% paper31-state-space.tex — State Space Exploration for Exhaustive Local Section Generation
% Paper 31 of the Judgment Geometry series.
% Compile: pdflatex -interaction=nonstopmode paper31-state-space.tex
\documentclass[11pt]{article}
\usepackage{lmodern}
% jugeo-common.tex — shared preamble for all Judgment Geometry papers
% Include via: % jugeo-common.tex — shared preamble for all Judgment Geometry papers
% Include via: % jugeo-common.tex — shared preamble for all Judgment Geometry papers
% Include via: \input{jugeo-common}

% ── Packages ────────────────────────────────────────────────────────────
\usepackage{amsmath,amssymb,amsthm,mathtools}
\usepackage{tikz-cd}
\usepackage{listings}
\usepackage{xcolor}
\usepackage{xspace}
\usepackage{booktabs}
\usepackage{multirow}
\usepackage{enumitem}
\usepackage[expansion=false]{microtype}
\usepackage{hyperref}
\usepackage{cleveref}
% \usepackage{thmtools}  % disabled: not available in this TeX installation
\usepackage{float}
\usepackage{caption}
\usepackage{subcaption}
\usepackage{tabularx}
\usepackage{stmaryrd}       % \llbracket, \rrbracket
\usepackage{bbm}            % \mathbbm{1}

% ── Page geometry ───────────────────────────────────────────────────────
\usepackage[margin=1in]{geometry}

% ── Theorem environments ────────────────────────────────────────────────
\theoremstyle{plain}
\newtheorem{theorem}{Theorem}[section]
\newtheorem{lemma}[theorem]{Lemma}
\newtheorem{proposition}[theorem]{Proposition}
\newtheorem{corollary}[theorem]{Corollary}

\theoremstyle{definition}
\newtheorem{definition}[theorem]{Definition}
\newtheorem{example}[theorem]{Example}
\newtheorem{construction}[theorem]{Construction}

\theoremstyle{remark}
\newtheorem{remark}[theorem]{Remark}
\newtheorem{notation}[theorem]{Notation}

% ── Python listings style ──────────────────────────────────────────────
\definecolor{jugeoBlue}{HTML}{2563EB}
\definecolor{jugeoGreen}{HTML}{059669}
\definecolor{jugeoRed}{HTML}{DC2626}
\definecolor{jugeoPurple}{HTML}{7C3AED}
\definecolor{jugeoGray}{HTML}{6B7280}
\definecolor{jugeoBg}{HTML}{F8FAFC}

\lstdefinestyle{jugeo-python}{
  language=Python,
  basicstyle=\ttfamily\small,
  keywordstyle=\color{jugeoBlue}\bfseries,
  stringstyle=\color{jugeoGreen},
  commentstyle=\color{jugeoGray}\itshape,
  numberstyle=\tiny\color{jugeoGray},
  backgroundcolor=\color{jugeoBg},
  numbers=left,
  numbersep=6pt,
  frame=single,
  framerule=0.4pt,
  rulecolor=\color{jugeoGray!40},
  breaklines=true,
  breakatwhitespace=true,
  showstringspaces=false,
  tabsize=4,
  xleftmargin=1.5em,
  framexleftmargin=1.5em,
  morekeywords={dataclass,frozen,field,Enum,IntEnum,auto,Any,
                Mapping,Sequence,Optional,match,case},
  literate={→}{{$\to$}}1 {←}{{$\leftarrow$}}1
           {≤}{{$\leq$}}1 {≥}{{$\geq$}}1 {≠}{{$\neq$}}1
           {∈}{{$\in$}}1 {∀}{{$\forall$}}1 {∃}{{$\exists$}}1
           {⊕}{{$\oplus$}}1 {⊖}{{$\ominus$}}1,
  escapeinside={(*@}{@*)},
}
\lstset{style=jugeo-python}

\lstdefinestyle{jugeo-output}{
  basicstyle=\ttfamily\small,
  backgroundcolor=\color{jugeoBg},
  frame=single,
  framerule=0.4pt,
  rulecolor=\color{jugeoGray!40},
  breaklines=true,
  showstringspaces=false,
  xleftmargin=1.5em,
  framexleftmargin=0.5em,
  numbers=none,
}

% ── JuGeo-specific macros ──────────────────────────────────────────────

% Judgment 8-tuple
\newcommand{\J}{\mathcal{J}}
\newcommand{\Jud}[8]{(#1,\,#2,\,#3,\,#4,\,#5,\,#6,\,#7,\,#8)}
\newcommand{\JudShort}[1]{\J_{#1}}

% Components
\newcommand{\coord}{\mathsf{c}}            % coordinate
\newcommand{\prop}{\varphi}                 % proposition
\newcommand{\carrier}{\mathcal{A}}          % carrier type
\newcommand{\evid}{\mathcal{E}}             % evidence bundle
\newcommand{\oblig}{\mathcal{O}}            % obligations
\newcommand{\obstr}{\mathcal{B}}            % obstructions
\newcommand{\trust}{\mathcal{T}}            % trust annotation
\newcommand{\prov}{\Pi}                     % provenance

% Trust algebra
\newcommand{\Talg}{\mathfrak{T}}            % trust ordered algebra
\newcommand{\Eadm}{\mathcal{E}_{\mathrm{adm}}}
\newcommand{\tleq}{\preceq}                 % trust ordering
\newcommand{\tjoin}{\oplus}                 % conservative join
\newcommand{\tatten}{\ominus}               % attenuation
\newcommand{\tpromote}{\uparrow_\pi}        % promotion
\newcommand{\tdemote}{\downarrow_\chi}       % demotion

% Trust tiers
\newcommand{\Tcontra}{\ensuremath{\mathsf{CONTRADICTED}}}
\newcommand{\Tunver}{\ensuremath{\mathsf{UNVERIFIED}}}
\newcommand{\Tcopilot}{\ensuremath{\mathsf{COPILOT}}}
\newcommand{\Toracle}{\ensuremath{\mathsf{ORACLE}}}
\newcommand{\Truntime}{\ensuremath{\mathsf{RUNTIME}}}
\newcommand{\Tsolver}{\ensuremath{\mathsf{SOLVER}}}
\newcommand{\Tproof}{\ensuremath{\mathsf{PROOF}}}

% Site and geometry
\newcommand{\Site}{\mathcal{S}}             % semantic site
\newcommand{\Cov}{\mathrm{Cov}}            % covering family
\newcommand{\Ob}{\mathrm{Ob}}              % objects
\newcommand{\Mor}{\mathrm{Mor}}            % morphisms
\newcommand{\res}{\mathrm{res}}            % restriction
\newcommand{\Cech}{\check{C}}              % Čech complex
\newcommand{\Hcoh}[1]{H^{#1}}             % cohomology group

% Descent
\newcommand{\descent}{\mathrm{desc}}
\newcommand{\glue}{\mathrm{glue}}
\newcommand{\overlap}[2]{#1 \cap #2}

% Semantic moves
\newcommand{\VERIFY}{\textsc{Verify}}
\newcommand{\CONSTRUCT}{\textsc{Construct}}
\newcommand{\NORMALIZE}{\textsc{Normalize}}
\newcommand{\GLUE}{\textsc{Glue}}
\newcommand{\RESTRICT}{\textsc{Restrict}}
\newcommand{\TRANSPORT}{\textsc{Transport}}
\newcommand{\REPAIR}{\textsc{Repair}}
\newcommand{\PROMOTE}{\textsc{Promote}}
\newcommand{\CHALLENGE}{\textsc{Challenge}}
\newcommand{\ESCALATE}{\textsc{Escalate}}

% Fragments
\newcommand{\QFLIA}{\texttt{QF\_LIA}}
\newcommand{\QFLRA}{\texttt{QF\_LRA}}
\newcommand{\QFBV}{\texttt{QF\_BV}}
\newcommand{\QFUF}{\texttt{QF\_UF}}

% Misc
\newcommand{\Python}{\textsc{Python}}
\newcommand{\JuGeo}{\textsc{JuGeo}}
\newcommand{\LEAN}{\textsc{Lean}\,4}
\newcommand{\Fstar}{F$^{\star}$}
\newcommand{\Coq}{\textsc{Coq}}
\newcommand{\Dafny}{\textsc{Dafny}}

% Common shortcuts
\newcommand{\ie}{i.e.\@\xspace}
\newcommand{\eg}{e.g.\@\xspace}
\newcommand{\cf}{cf.\@\xspace}
\newcommand{\wrt}{w.r.t.\@\xspace}

% Evaluation shortcuts
\newcommand{\numcases}{300}
\newcommand{\numspec}{100}
\newcommand{\numeq}{100}
\newcommand{\numbug}{100}
\newcommand{\medtime}{6.5\,ms}
\newcommand{\totaltime}{1.949\,s}


% ── Packages ────────────────────────────────────────────────────────────
\usepackage{amsmath,amssymb,amsthm,mathtools}
\usepackage{tikz-cd}
\usepackage{listings}
\usepackage{xcolor}
\usepackage{xspace}
\usepackage{booktabs}
\usepackage{multirow}
\usepackage{enumitem}
\usepackage[expansion=false]{microtype}
\usepackage{hyperref}
\usepackage{cleveref}
% \usepackage{thmtools}  % disabled: not available in this TeX installation
\usepackage{float}
\usepackage{caption}
\usepackage{subcaption}
\usepackage{tabularx}
\usepackage{stmaryrd}       % \llbracket, \rrbracket
\usepackage{bbm}            % \mathbbm{1}

% ── Page geometry ───────────────────────────────────────────────────────
\usepackage[margin=1in]{geometry}

% ── Theorem environments ────────────────────────────────────────────────
\theoremstyle{plain}
\newtheorem{theorem}{Theorem}[section]
\newtheorem{lemma}[theorem]{Lemma}
\newtheorem{proposition}[theorem]{Proposition}
\newtheorem{corollary}[theorem]{Corollary}

\theoremstyle{definition}
\newtheorem{definition}[theorem]{Definition}
\newtheorem{example}[theorem]{Example}
\newtheorem{construction}[theorem]{Construction}

\theoremstyle{remark}
\newtheorem{remark}[theorem]{Remark}
\newtheorem{notation}[theorem]{Notation}

% ── Python listings style ──────────────────────────────────────────────
\definecolor{jugeoBlue}{HTML}{2563EB}
\definecolor{jugeoGreen}{HTML}{059669}
\definecolor{jugeoRed}{HTML}{DC2626}
\definecolor{jugeoPurple}{HTML}{7C3AED}
\definecolor{jugeoGray}{HTML}{6B7280}
\definecolor{jugeoBg}{HTML}{F8FAFC}

\lstdefinestyle{jugeo-python}{
  language=Python,
  basicstyle=\ttfamily\small,
  keywordstyle=\color{jugeoBlue}\bfseries,
  stringstyle=\color{jugeoGreen},
  commentstyle=\color{jugeoGray}\itshape,
  numberstyle=\tiny\color{jugeoGray},
  backgroundcolor=\color{jugeoBg},
  numbers=left,
  numbersep=6pt,
  frame=single,
  framerule=0.4pt,
  rulecolor=\color{jugeoGray!40},
  breaklines=true,
  breakatwhitespace=true,
  showstringspaces=false,
  tabsize=4,
  xleftmargin=1.5em,
  framexleftmargin=1.5em,
  morekeywords={dataclass,frozen,field,Enum,IntEnum,auto,Any,
                Mapping,Sequence,Optional,match,case},
  literate={→}{{$\to$}}1 {←}{{$\leftarrow$}}1
           {≤}{{$\leq$}}1 {≥}{{$\geq$}}1 {≠}{{$\neq$}}1
           {∈}{{$\in$}}1 {∀}{{$\forall$}}1 {∃}{{$\exists$}}1
           {⊕}{{$\oplus$}}1 {⊖}{{$\ominus$}}1,
  escapeinside={(*@}{@*)},
}
\lstset{style=jugeo-python}

\lstdefinestyle{jugeo-output}{
  basicstyle=\ttfamily\small,
  backgroundcolor=\color{jugeoBg},
  frame=single,
  framerule=0.4pt,
  rulecolor=\color{jugeoGray!40},
  breaklines=true,
  showstringspaces=false,
  xleftmargin=1.5em,
  framexleftmargin=0.5em,
  numbers=none,
}

% ── JuGeo-specific macros ──────────────────────────────────────────────

% Judgment 8-tuple
\newcommand{\J}{\mathcal{J}}
\newcommand{\Jud}[8]{(#1,\,#2,\,#3,\,#4,\,#5,\,#6,\,#7,\,#8)}
\newcommand{\JudShort}[1]{\J_{#1}}

% Components
\newcommand{\coord}{\mathsf{c}}            % coordinate
\newcommand{\prop}{\varphi}                 % proposition
\newcommand{\carrier}{\mathcal{A}}          % carrier type
\newcommand{\evid}{\mathcal{E}}             % evidence bundle
\newcommand{\oblig}{\mathcal{O}}            % obligations
\newcommand{\obstr}{\mathcal{B}}            % obstructions
\newcommand{\trust}{\mathcal{T}}            % trust annotation
\newcommand{\prov}{\Pi}                     % provenance

% Trust algebra
\newcommand{\Talg}{\mathfrak{T}}            % trust ordered algebra
\newcommand{\Eadm}{\mathcal{E}_{\mathrm{adm}}}
\newcommand{\tleq}{\preceq}                 % trust ordering
\newcommand{\tjoin}{\oplus}                 % conservative join
\newcommand{\tatten}{\ominus}               % attenuation
\newcommand{\tpromote}{\uparrow_\pi}        % promotion
\newcommand{\tdemote}{\downarrow_\chi}       % demotion

% Trust tiers
\newcommand{\Tcontra}{\ensuremath{\mathsf{CONTRADICTED}}}
\newcommand{\Tunver}{\ensuremath{\mathsf{UNVERIFIED}}}
\newcommand{\Tcopilot}{\ensuremath{\mathsf{COPILOT}}}
\newcommand{\Toracle}{\ensuremath{\mathsf{ORACLE}}}
\newcommand{\Truntime}{\ensuremath{\mathsf{RUNTIME}}}
\newcommand{\Tsolver}{\ensuremath{\mathsf{SOLVER}}}
\newcommand{\Tproof}{\ensuremath{\mathsf{PROOF}}}

% Site and geometry
\newcommand{\Site}{\mathcal{S}}             % semantic site
\newcommand{\Cov}{\mathrm{Cov}}            % covering family
\newcommand{\Ob}{\mathrm{Ob}}              % objects
\newcommand{\Mor}{\mathrm{Mor}}            % morphisms
\newcommand{\res}{\mathrm{res}}            % restriction
\newcommand{\Cech}{\check{C}}              % Čech complex
\newcommand{\Hcoh}[1]{H^{#1}}             % cohomology group

% Descent
\newcommand{\descent}{\mathrm{desc}}
\newcommand{\glue}{\mathrm{glue}}
\newcommand{\overlap}[2]{#1 \cap #2}

% Semantic moves
\newcommand{\VERIFY}{\textsc{Verify}}
\newcommand{\CONSTRUCT}{\textsc{Construct}}
\newcommand{\NORMALIZE}{\textsc{Normalize}}
\newcommand{\GLUE}{\textsc{Glue}}
\newcommand{\RESTRICT}{\textsc{Restrict}}
\newcommand{\TRANSPORT}{\textsc{Transport}}
\newcommand{\REPAIR}{\textsc{Repair}}
\newcommand{\PROMOTE}{\textsc{Promote}}
\newcommand{\CHALLENGE}{\textsc{Challenge}}
\newcommand{\ESCALATE}{\textsc{Escalate}}

% Fragments
\newcommand{\QFLIA}{\texttt{QF\_LIA}}
\newcommand{\QFLRA}{\texttt{QF\_LRA}}
\newcommand{\QFBV}{\texttt{QF\_BV}}
\newcommand{\QFUF}{\texttt{QF\_UF}}

% Misc
\newcommand{\Python}{\textsc{Python}}
\newcommand{\JuGeo}{\textsc{JuGeo}}
\newcommand{\LEAN}{\textsc{Lean}\,4}
\newcommand{\Fstar}{F$^{\star}$}
\newcommand{\Coq}{\textsc{Coq}}
\newcommand{\Dafny}{\textsc{Dafny}}

% Common shortcuts
\newcommand{\ie}{i.e.\@\xspace}
\newcommand{\eg}{e.g.\@\xspace}
\newcommand{\cf}{cf.\@\xspace}
\newcommand{\wrt}{w.r.t.\@\xspace}

% Evaluation shortcuts
\newcommand{\numcases}{300}
\newcommand{\numspec}{100}
\newcommand{\numeq}{100}
\newcommand{\numbug}{100}
\newcommand{\medtime}{6.5\,ms}
\newcommand{\totaltime}{1.949\,s}


% ── Packages ────────────────────────────────────────────────────────────
\usepackage{amsmath,amssymb,amsthm,mathtools}
\usepackage{tikz-cd}
\usepackage{listings}
\usepackage{xcolor}
\usepackage{xspace}
\usepackage{booktabs}
\usepackage{multirow}
\usepackage{enumitem}
\usepackage[expansion=false]{microtype}
\usepackage{hyperref}
\usepackage{cleveref}
% \usepackage{thmtools}  % disabled: not available in this TeX installation
\usepackage{float}
\usepackage{caption}
\usepackage{subcaption}
\usepackage{tabularx}
\usepackage{stmaryrd}       % \llbracket, \rrbracket
\usepackage{bbm}            % \mathbbm{1}

% ── Page geometry ───────────────────────────────────────────────────────
\usepackage[margin=1in]{geometry}

% ── Theorem environments ────────────────────────────────────────────────
\theoremstyle{plain}
\newtheorem{theorem}{Theorem}[section]
\newtheorem{lemma}[theorem]{Lemma}
\newtheorem{proposition}[theorem]{Proposition}
\newtheorem{corollary}[theorem]{Corollary}

\theoremstyle{definition}
\newtheorem{definition}[theorem]{Definition}
\newtheorem{example}[theorem]{Example}
\newtheorem{construction}[theorem]{Construction}

\theoremstyle{remark}
\newtheorem{remark}[theorem]{Remark}
\newtheorem{notation}[theorem]{Notation}

% ── Python listings style ──────────────────────────────────────────────
\definecolor{jugeoBlue}{HTML}{2563EB}
\definecolor{jugeoGreen}{HTML}{059669}
\definecolor{jugeoRed}{HTML}{DC2626}
\definecolor{jugeoPurple}{HTML}{7C3AED}
\definecolor{jugeoGray}{HTML}{6B7280}
\definecolor{jugeoBg}{HTML}{F8FAFC}

\lstdefinestyle{jugeo-python}{
  language=Python,
  basicstyle=\ttfamily\small,
  keywordstyle=\color{jugeoBlue}\bfseries,
  stringstyle=\color{jugeoGreen},
  commentstyle=\color{jugeoGray}\itshape,
  numberstyle=\tiny\color{jugeoGray},
  backgroundcolor=\color{jugeoBg},
  numbers=left,
  numbersep=6pt,
  frame=single,
  framerule=0.4pt,
  rulecolor=\color{jugeoGray!40},
  breaklines=true,
  breakatwhitespace=true,
  showstringspaces=false,
  tabsize=4,
  xleftmargin=1.5em,
  framexleftmargin=1.5em,
  morekeywords={dataclass,frozen,field,Enum,IntEnum,auto,Any,
                Mapping,Sequence,Optional,match,case},
  literate={→}{{$\to$}}1 {←}{{$\leftarrow$}}1
           {≤}{{$\leq$}}1 {≥}{{$\geq$}}1 {≠}{{$\neq$}}1
           {∈}{{$\in$}}1 {∀}{{$\forall$}}1 {∃}{{$\exists$}}1
           {⊕}{{$\oplus$}}1 {⊖}{{$\ominus$}}1,
  escapeinside={(*@}{@*)},
}
\lstset{style=jugeo-python}

\lstdefinestyle{jugeo-output}{
  basicstyle=\ttfamily\small,
  backgroundcolor=\color{jugeoBg},
  frame=single,
  framerule=0.4pt,
  rulecolor=\color{jugeoGray!40},
  breaklines=true,
  showstringspaces=false,
  xleftmargin=1.5em,
  framexleftmargin=0.5em,
  numbers=none,
}

% ── JuGeo-specific macros ──────────────────────────────────────────────

% Judgment 8-tuple
\newcommand{\J}{\mathcal{J}}
\newcommand{\Jud}[8]{(#1,\,#2,\,#3,\,#4,\,#5,\,#6,\,#7,\,#8)}
\newcommand{\JudShort}[1]{\J_{#1}}

% Components
\newcommand{\coord}{\mathsf{c}}            % coordinate
\newcommand{\prop}{\varphi}                 % proposition
\newcommand{\carrier}{\mathcal{A}}          % carrier type
\newcommand{\evid}{\mathcal{E}}             % evidence bundle
\newcommand{\oblig}{\mathcal{O}}            % obligations
\newcommand{\obstr}{\mathcal{B}}            % obstructions
\newcommand{\trust}{\mathcal{T}}            % trust annotation
\newcommand{\prov}{\Pi}                     % provenance

% Trust algebra
\newcommand{\Talg}{\mathfrak{T}}            % trust ordered algebra
\newcommand{\Eadm}{\mathcal{E}_{\mathrm{adm}}}
\newcommand{\tleq}{\preceq}                 % trust ordering
\newcommand{\tjoin}{\oplus}                 % conservative join
\newcommand{\tatten}{\ominus}               % attenuation
\newcommand{\tpromote}{\uparrow_\pi}        % promotion
\newcommand{\tdemote}{\downarrow_\chi}       % demotion

% Trust tiers
\newcommand{\Tcontra}{\ensuremath{\mathsf{CONTRADICTED}}}
\newcommand{\Tunver}{\ensuremath{\mathsf{UNVERIFIED}}}
\newcommand{\Tcopilot}{\ensuremath{\mathsf{COPILOT}}}
\newcommand{\Toracle}{\ensuremath{\mathsf{ORACLE}}}
\newcommand{\Truntime}{\ensuremath{\mathsf{RUNTIME}}}
\newcommand{\Tsolver}{\ensuremath{\mathsf{SOLVER}}}
\newcommand{\Tproof}{\ensuremath{\mathsf{PROOF}}}

% Site and geometry
\newcommand{\Site}{\mathcal{S}}             % semantic site
\newcommand{\Cov}{\mathrm{Cov}}            % covering family
\newcommand{\Ob}{\mathrm{Ob}}              % objects
\newcommand{\Mor}{\mathrm{Mor}}            % morphisms
\newcommand{\res}{\mathrm{res}}            % restriction
\newcommand{\Cech}{\check{C}}              % Čech complex
\newcommand{\Hcoh}[1]{H^{#1}}             % cohomology group

% Descent
\newcommand{\descent}{\mathrm{desc}}
\newcommand{\glue}{\mathrm{glue}}
\newcommand{\overlap}[2]{#1 \cap #2}

% Semantic moves
\newcommand{\VERIFY}{\textsc{Verify}}
\newcommand{\CONSTRUCT}{\textsc{Construct}}
\newcommand{\NORMALIZE}{\textsc{Normalize}}
\newcommand{\GLUE}{\textsc{Glue}}
\newcommand{\RESTRICT}{\textsc{Restrict}}
\newcommand{\TRANSPORT}{\textsc{Transport}}
\newcommand{\REPAIR}{\textsc{Repair}}
\newcommand{\PROMOTE}{\textsc{Promote}}
\newcommand{\CHALLENGE}{\textsc{Challenge}}
\newcommand{\ESCALATE}{\textsc{Escalate}}

% Fragments
\newcommand{\QFLIA}{\texttt{QF\_LIA}}
\newcommand{\QFLRA}{\texttt{QF\_LRA}}
\newcommand{\QFBV}{\texttt{QF\_BV}}
\newcommand{\QFUF}{\texttt{QF\_UF}}

% Misc
\newcommand{\Python}{\textsc{Python}}
\newcommand{\JuGeo}{\textsc{JuGeo}}
\newcommand{\LEAN}{\textsc{Lean}\,4}
\newcommand{\Fstar}{F$^{\star}$}
\newcommand{\Coq}{\textsc{Coq}}
\newcommand{\Dafny}{\textsc{Dafny}}

% Common shortcuts
\newcommand{\ie}{i.e.\@\xspace}
\newcommand{\eg}{e.g.\@\xspace}
\newcommand{\cf}{cf.\@\xspace}
\newcommand{\wrt}{w.r.t.\@\xspace}

% Evaluation shortcuts
\newcommand{\numcases}{300}
\newcommand{\numspec}{100}
\newcommand{\numeq}{100}
\newcommand{\numbug}{100}
\newcommand{\medtime}{6.5\,ms}
\newcommand{\totaltime}{1.949\,s}

% experiment-data.tex — AUTO-GENERATED from real jugeo CLI runs
% DO NOT EDIT — regenerate with: python3 experiments/run_universal_metrics.py
% Generated from 30 programs across 6 domains

\newcommand{\expTotalPrograms}{30}
\newcommand{\expVerified}{30}
\newcommand{\expAccuracy}{100.0\%}
\newcommand{\expCoordsMin}{2}
\newcommand{\expCoordsMax}{10}
\newcommand{\expCoordsMean}{5.2}
\newcommand{\expCoordsMedian}{5.5}
\newcommand{\expPropsMin}{4}
\newcommand{\expPropsMax}{20}
\newcommand{\expPropsMean}{10.4}
\newcommand{\expPropsSum}{311}
\newcommand{\expPropsOkSum}{310}
\newcommand{\expTimeMin}{3.506\,s}
\newcommand{\expTimeMax}{6.714\,s}
\newcommand{\expTimeMean}{4.64\,s}
\newcommand{\expTimeMedian}{4.508\,s}
\newcommand{\expTimeTotal}{139.204\,s}
\newcommand{\expObstructionTotal}{0}
\newcommand{\expHOneZero}{30}
\newcommand{\expDomainCount}{6}
\newcommand{\expTrustCOPILOTSUGGESTED}{30}
\newcommand{\expDomainDsCount}{5}
\newcommand{\expDomainDsVerified}{5}
\newcommand{\expDomainDsAvgCoords}{7.6}
\newcommand{\expDomainDsAvgProps}{15.2}
\newcommand{\expDomainMathCount}{5}
\newcommand{\expDomainMathVerified}{5}
\newcommand{\expDomainMathAvgCoords}{4.8}
\newcommand{\expDomainMathAvgProps}{9.4}
\newcommand{\expDomainSmCount}{5}
\newcommand{\expDomainSmVerified}{5}
\newcommand{\expDomainSmAvgCoords}{7.6}
\newcommand{\expDomainSmAvgProps}{15.2}
\newcommand{\expDomainSortCount}{5}
\newcommand{\expDomainSortVerified}{5}
\newcommand{\expDomainSortAvgCoords}{2.4}
\newcommand{\expDomainSortAvgProps}{4.8}
\newcommand{\expDomainStrCount}{5}
\newcommand{\expDomainStrVerified}{5}
\newcommand{\expDomainStrAvgCoords}{3.2}
\newcommand{\expDomainStrAvgProps}{6.4}
\newcommand{\expDomainWebCount}{5}
\newcommand{\expDomainWebVerified}{5}
\newcommand{\expDomainWebAvgCoords}{5.6}
\newcommand{\expDomainWebAvgProps}{11.2}

% subsystem-data.tex — AUTO-GENERATED from real jugeo subsystem experiments
% DO NOT EDIT — regenerate with: python3 experiments/run_subsystem_experiments.py

\newcommand{\subCliTotal}{10}
\newcommand{\subCliVerified}{9}
\newcommand{\subCliAccuracy}{90.0\%}
\newcommand{\subCliMeanTime}{4.132\,s}
\newcommand{\subCliMinTime}{3.472\,s}
\newcommand{\subCliMaxTime}{5.372\,s}
\newcommand{\subCliTotalCoords}{54}
\newcommand{\subCliTotalProps}{107}
\newcommand{\subCliTotalPropsOk}{106}
\newcommand{\subSiteCalls}{90}
\newcommand{\subSiteSuccessful}{69}
\newcommand{\subSiteSuccessRate}{76.7\%}
\newcommand{\subSiteMeanTime}{0.0001\,s}
\newcommand{\subSiteTotalTime}{0.005\,s}
\newcommand{\subSpecificationsatisfactionSmallTime}{0.0003\,s}
\newcommand{\subSpecificationsatisfactionMediumTime}{0.0001\,s}
\newcommand{\subSpecificationsatisfactionLargeTime}{0.0001\,s}
\newcommand{\subInhabitantfleetSmallTime}{0.0\,s}
\newcommand{\subInhabitantfleetMediumTime}{0.0\,s}
\newcommand{\subInhabitantfleetLargeTime}{0.0\,s}
\newcommand{\subTheoremecologySmallTime}{0.0\,s}
\newcommand{\subTheoremecologyMediumTime}{0.0\,s}
\newcommand{\subTheoremecologyLargeTime}{0.0\,s}
\newcommand{\subStatespaceexplorationSmallTime}{0.0\,s}
\newcommand{\subStatespaceexplorationMediumTime}{0.0\,s}
\newcommand{\subStatespaceexplorationLargeTime}{0.0\,s}
\newcommand{\subRepairsemanticsSmallTime}{0.0\,s}
\newcommand{\subRepairsemanticsMediumTime}{0.0\,s}
\newcommand{\subRepairsemanticsLargeTime}{0.0\,s}
\newcommand{\subReplaygluingSmallTime}{0.0\,s}
\newcommand{\subReplaygluingMediumTime}{0.0\,s}
\newcommand{\subReplaygluingLargeTime}{0.0\,s}
\newcommand{\subSemanticfuturesSmallTime}{0.0\,s}
\newcommand{\subSemanticfuturesMediumTime}{0.0\,s}
\newcommand{\subSemanticfuturesLargeTime}{0.0\,s}
\newcommand{\subHypercovertreatySmallTime}{0.0\,s}
\newcommand{\subHypercovertreatyMediumTime}{0.0\,s}
\newcommand{\subHypercovertreatyLargeTime}{0.0\,s}
\newcommand{\subEncodeforsolverSmallTime}{0.0026\,s}
\newcommand{\subEncodeforsolverMediumTime}{0.0002\,s}
\newcommand{\subEncodeforsolverLargeTime}{0.0001\,s}
\newcommand{\subInterfaceroutingSmallTime}{0.0\,s}
\newcommand{\subInterfaceroutingMediumTime}{0.0\,s}
\newcommand{\subInterfaceroutingLargeTime}{0.0\,s}
\newcommand{\subOrchestrateverificationSmallTime}{0.0002\,s}
\newcommand{\subOrchestrateverificationMediumTime}{0.0\,s}
\newcommand{\subOrchestrateverificationLargeTime}{0.0\,s}
\newcommand{\subRunfulldescentSmallTime}{0.0\,s}
\newcommand{\subRunfulldescentMediumTime}{0.0\,s}
\newcommand{\subRunfulldescentLargeTime}{0.0\,s}
\newcommand{\subKernellifecycleSmallTime}{0.0\,s}
\newcommand{\subKernellifecycleMediumTime}{0.0\,s}
\newcommand{\subKernellifecycleLargeTime}{0.0\,s}
\newcommand{\subDiscoverypipelineSmallTime}{0.0\,s}
\newcommand{\subDiscoverypipelineMediumTime}{0.0\,s}
\newcommand{\subDiscoverypipelineLargeTime}{0.0\,s}
\newcommand{\subAnalogytransportSmallTime}{0.0\,s}
\newcommand{\subAnalogytransportMediumTime}{0.0\,s}
\newcommand{\subAnalogytransportLargeTime}{0.0\,s}
\newcommand{\subFormalcoresiteSmallTime}{0.0001\,s}
\newcommand{\subFormalcoresiteMediumTime}{0.0\,s}
\newcommand{\subFormalcoresiteLargeTime}{0.0\,s}
\newcommand{\subProblematlasSmallTime}{0.0\,s}
\newcommand{\subProblematlasMediumTime}{0.0\,s}
\newcommand{\subProblematlasLargeTime}{0.0\,s}
\newcommand{\subPublicalignmentSmallTime}{0.0\,s}
\newcommand{\subPublicalignmentMediumTime}{0.0\,s}
\newcommand{\subPublicalignmentLargeTime}{0.0\,s}
\newcommand{\subRegimebootstrappingSmallTime}{0.0\,s}
\newcommand{\subRegimebootstrappingMediumTime}{0.0\,s}
\newcommand{\subRegimebootstrappingLargeTime}{0.0\,s}
\newcommand{\subRelationalrefinementSmallTime}{0.0\,s}
\newcommand{\subRelationalrefinementMediumTime}{0.0\,s}
\newcommand{\subRelationalrefinementLargeTime}{0.0\,s}
\newcommand{\subSemanticclosureSmallTime}{0.0\,s}
\newcommand{\subSemanticclosureMediumTime}{0.0\,s}
\newcommand{\subSemanticclosureLargeTime}{0.0\,s}
\newcommand{\subTheoremeconomicsSmallTime}{0.0001\,s}
\newcommand{\subTheoremeconomicsMediumTime}{0.0\,s}
\newcommand{\subTheoremeconomicsLargeTime}{0.0\,s}
\newcommand{\subBenchmarksuiteSmallTime}{0.0\,s}
\newcommand{\subBenchmarksuiteMediumTime}{0.0\,s}
\newcommand{\subBenchmarksuiteLargeTime}{0.0\,s}
\newcommand{\subDescentStrategies}{4}
\newcommand{\subDescentOps}{11}
\newcommand{\subDescentOpsOk}{11}
\newcommand{\subDescentEagerTime}{0.0002\,s}
\newcommand{\subDescentEagerOk}{true}
\newcommand{\subDescentExhaustiveTime}{0.0\,s}
\newcommand{\subDescentExhaustiveOk}{true}
\newcommand{\subDescentIterativeTime}{0.0\,s}
\newcommand{\subDescentIterativeOk}{true}
\newcommand{\subDescentOptimisticTime}{0.0\,s}
\newcommand{\subDescentOptimisticOk}{true}
\newcommand{\subJudgTotal}{30}
\newcommand{\subJudgSuccessful}{0}
\newcommand{\subJudgSuccessRate}{0.0\%}
\newcommand{\subJudgMeanTimeUs}{7.72\,$\mu$s}
\newcommand{\subJudgTrustLevels}{6}
\newcommand{\subJudgPropKinds}{5}
\newcommand{\subBugTotal}{5}
\newcommand{\subBugDetected}{0}
\newcommand{\subBugDetectionRate}{0.0\%}
\newcommand{\subBugFalsePositiveRate}{0.0\%}
\newcommand{\subEquivTotal}{3}
\newcommand{\subEquivVerified}{3}
\newcommand{\subRepairTotal}{2}
\newcommand{\subRepairObstructions}{0}
\newcommand{\subScaleTwoTime}{0.0001\,s}
\newcommand{\subScaleFourTime}{0.0\,s}
\newcommand{\subScaleEightTime}{0.0\,s}
\newcommand{\subScaleTwelveTime}{0.0\,s}
\newcommand{\subScaleSixteenTime}{0.0\,s}
\newcommand{\subScaleTwentyTime}{0.0\,s}
\newcommand{\subTotalExperimentTime}{95.6\,s}

% data-paper31.tex — AUTO-GENERATED by exp31_state_space.py
% DO NOT EDIT — regenerate with: python3 experiments/exp31_state_space.py

\newcommand{\ppXXXIBfsStates}{16}
\newcommand{\ppXXXIBfsTimeMs}{3928.8\,ms}
\newcommand{\ppXXXIBfsCoverage}{1.00}
\newcommand{\ppXXXIBfsRounds}{3.1}
\newcommand{\ppXXXIDfsStates}{17}
\newcommand{\ppXXXIDfsTimeMs}{4911.0\,ms}
\newcommand{\ppXXXIDfsCoverage}{1.00}
\newcommand{\ppXXXIDfsRounds}{3.6}
\newcommand{\ppXXXIBmcStates}{9}
\newcommand{\ppXXXIBmcTimeMs}{2750.2\,ms}
\newcommand{\ppXXXIBmcCoverage}{0.98}
\newcommand{\ppXXXIBmcRounds}{1.9}
\newcommand{\ppXXXIHeurStates}{13}
\newcommand{\ppXXXIHeurTimeMs}{3536.0\,ms}
\newcommand{\ppXXXIHeurCoverage}{1.00}
\newcommand{\ppXXXIHeurRounds}{3.1}
\newcommand{\ppXXXISmallBfsMs}{3650.1\,ms}
\newcommand{\ppXXXISmallDfsMs}{4197.6\,ms}
\newcommand{\ppXXXIMedBfsMs}{3712.5\,ms}
\newcommand{\ppXXXIMedDfsMs}{4640.6\,ms}
\newcommand{\ppXXXILargeBfsMs}{6250.3\,ms}
\newcommand{\ppXXXILargeDfsMs}{8750.4\,ms}
\newcommand{\ppXXXIXlargeBfsMs}{0.0\,ms}
\newcommand{\ppXXXIXlargeDfsMs}{0.0\,ms}
\newcommand{\ppXXXITotalPrograms}{10}


% ── Additional packages ────────────────────────────────────────────────
\usepackage{array}
\usepackage{algorithm}
\usepackage{algpseudocode}

% ── Paper-specific macros ──────────────────────────────────────────────
\newcommand{\St}{\Sigma}                      % state space
\newcommand{\state}{\sigma}                   % individual state
\newcommand{\frontier}{\mathcal{F}}           % frontier / worklist
\newcommand{\visited}{\mathcal{V}}            % visited set
\newcommand{\Trans}{\delta}                   % transition relation
\newcommand{\reach}{\mathrm{Reach}}           % reachable states
\newcommand{\BFS}{\textsc{Bfs}}
\newcommand{\DFS}{\textsc{Dfs}}
\newcommand{\BMC}{\textsc{Bmc}}
\newcommand{\HEUR}{\textsc{Heur}}
\newcommand{\depth}{\mathsf{d}}               % depth bound
\newcommand{\heur}{h}                         % heuristic function
\newcommand{\LocalSec}{\ell}                  % local section
\newcommand{\LocSecs}{\mathcal{L}}            % set of local sections
\newcommand{\LCon}{\mathsf{LC}}               % LocalConstructor
\newcommand{\Orch}{\mathsf{Orch}}             % ConstructionOrchestrator
\newcommand{\CStatus}{\mathsf{CS}}            % ConstructionStatus
\newcommand{\CCond}{\mathsf{CC}}              % CompletionCondition
\newcommand{\SearchStrat}{\mathsf{SS}}        % SearchStrategy
\newcommand{\StateExp}{\mathsf{SE}}           % StateExplorer
\newcommand{\patchset}{\mathcal{P}}           % patch set
\renewcommand{\coord}{\mathsf{c}}               % coordinate (already in common)
\newcommand{\budg}{B}                         % budget
\newcommand{\covfrac}{f}                      % coverage fraction
\newcommand{\Treach}{T_{\mathrm{reach}}}      % reachable-state time

\title{\textbf{State Space Exploration for Exhaustive\\Local Section Generation}}

\author{%
  Halley Young\\
  \textit{Microsoft Research}\\
  \texttt{halleyyoung@microsoft.com}
}

\date{\today}

\begin{document}
\maketitle

% =====================================================================
\begin{abstract}
We present a \emph{state space exploration} framework for the
systematic, exhaustive generation of local sections in the Judgment
Geometry (\JuGeo{}) program-verification system.  A \emph{state space}
$(\St, \Trans, \state_0)$ models a program's verification process as a
directed graph whose nodes are semantic states---partial assignments of
sections to coordinate patches---and whose edges are the six canonical
semantic moves.  A \emph{state explorer} (\texttt{StateExplorer})
traverses this graph using a pluggable \texttt{SearchStrategy}, which
can be breadth-first (\BFS), depth-first (\DFS), bounded model checking
(\BMC), or heuristically guided (\HEUR).  The explorer passes each
reached state to a \texttt{LocalConstructor}, which builds the local
section at the corresponding coordinate; the \texttt{ConstructionOrchestrator}
mediates across multiple coordinates and tracks \texttt{ConstructionStatus}.
Exploration terminates according to a \texttt{CompletionCondition} that
abstracts away the concrete strategy.

Our main theoretical contribution is the \emph{Bounded Completeness
Theorem}: for any program whose reachable state space has at most $N$
states, \BFS{} exploration visits every reachable state exactly once
and delivers a complete set of local sections in time $O(N + |\Trans|)$.
Completeness is measured via a novel \emph{coverage fraction} $\covfrac$
that counts the proportion of coordinate patches for which a valid local
section has been accepted.  We prove that \BFS{} drives $\covfrac$ to
$1$ for finite-state programs, while \DFS{} is complete only modulo a
cycle-detection invariant.  Bounded model checking provides a tunable
trade-off: with depth bound $\depth$, the fraction of covered patches
satisfies $\covfrac \ge 1 - (1 - p)^{\depth}$ where $p$ is the
single-step coverage probability.

We integrate the framework into \JuGeo{} and connect it to the
local-construction subsystem in
\texttt{src/jugeo/generation/local\_construction/}.  The repository's
measured benchmark covers 10 programs: \BFS{} visits 16 states on
average with 3928.8\,ms mean time and full coverage, \DFS{} visits 17
states with 4911.0\,ms mean time and full coverage, \BMC{} visits fewer
states (9 on average) with 0.98 coverage, and the heuristic strategy
matches full coverage at 3536.0\,ms mean time.
\end{abstract}

\newpage

% =====================================================================
\section{Introduction}\label{sec:intro}
% =====================================================================

Program verification in \JuGeo{} is organised around a
\emph{sheaf-theoretic} model: a program is decomposed into a cover of
coordinate patches $\patchset = \{p_1, \dots, p_n\}$, and verification
proceeds by constructing a \emph{local section} $\LocalSec_{p_i}$ at
each patch.  Global soundness is then ensured by a descent check: the
local sections glue consistently into a global section.

\paragraph{The coverage problem.}
In practice, not all coordinate patches are obvious from the program
text.  Control flow, data-dependent branching, and exception paths
give rise to \emph{latent coordinates} that are only exposed when the
program is executed or symbolically evaluated.  Discovering all latent
coordinates---and building local sections for each one---is the
\emph{coverage problem} for verification generation.

\paragraph{State spaces as the natural model.}
A state space $(\St, \Trans, \state_0)$ is the natural model for this
discovery process.  Each state $\state \in \St$ is a snapshot of the
current generation process: which patches have sections, which
obligations are open, which treaties are active.  A transition
$\Trans(\state, \state')$ fires when a semantic move advances
the process.  A successful state is one in which all patches have
accepted sections.

\paragraph{The exploration gap.}
Prior work in \JuGeo{} (Papers~1--30) has focused on individual
semantic moves and their algebraic properties.  The problem of
\emph{systematically scheduling} those moves across the state space has
not been studied.  Two deficiencies arise in an ad-hoc scheduler:
\begin{enumerate}[label=(\roman*)]
  \item \emph{Redundancy}: the same state may be reached by multiple
    move sequences and reconstructed from scratch each time.
  \item \emph{Incompleteness}: some reachable coordinates are never
    visited because the scheduler follows a single path.
\end{enumerate}
A principled state-space explorer eliminates both deficiencies.

\paragraph{Contributions.}
This paper makes four main contributions:
\begin{enumerate}
  \item \textbf{State space model} (\S\ref{sec:statespace}): a formal
    definition of \texttt{StateSpace} and \texttt{StateExplorer}, with
    a precise characterisation of reachability.
  \item \textbf{Search strategies} (\S\ref{sec:strategies}): the
    \texttt{SearchStrategy} enumeration---\BFS{}, \DFS{}, \BMC{},
    \HEUR{}---with correctness and complexity analyses.
  \item \textbf{Local construction integration}
    (\S\S\ref{sec:local}--\ref{sec:orchestration}): the connection from
    states to local sections via \texttt{LocalConstructor} and
    \texttt{ConstructionOrchestrator}.
  \item \textbf{Bounded completeness theorem}
    (\S\ref{sec:theorem}): a proof that \BFS{} is complete for
    finite-state programs, formalised in \LEAN{}.
\end{enumerate}

% =====================================================================
\section{State Space Model}\label{sec:statespace}
% =====================================================================

\subsection{Semantic states}

\begin{definition}[Semantic state]\label{def:state}
Let $\patchset = \{p_1, \dots, p_n\}$ be the patch cover of a program
$P$, and let $\mathcal{S}$ be the universe of available local sections.
A \emph{semantic state} is a triple
\[
  \state = (\alpha,\, \Omega,\, \mathcal{T})
\]
where $\alpha\colon \patchset \rightharpoonup \mathcal{S}$ is a partial
assignment of sections to patches, $\Omega$ is the set of open
obligations, and $\mathcal{T}$ is the set of active treaties.  We write
$\St$ for the set of all semantic states.
\end{definition}

A state $\state$ is \emph{initial} if $\alpha = \emptyset$,
$\Omega = \Omega_0(P)$ (the base obligations of $P$), and
$\mathcal{T} = \emptyset$.  A state is \emph{terminal} (\emph{goal})
if $\alpha$ is total (\ie{} all patches have sections), all
obligations in $\Omega$ are closed, and all treaties in $\mathcal{T}$
are ratified.

\begin{definition}[Transition relation]\label{def:trans}
The \emph{transition relation} $\Trans \subseteq \St \times \mathcal{M}
\times \St$ has one rule for each of the six canonical semantic moves:
\[
  \CONSTRUCT,\quad \VERIFY,\quad \NORMALIZE,\quad
  \GLUE,\quad \RESTRICT,\quad \REPAIR.
\]
We write $\state \xrightarrow{m} \state'$ when move $m$ transforms
$\state$ into $\state'$.  Moves are \emph{deterministic} given the
current state and a choice of patch/section: two distinct moves applied
to the same state always yield distinct successors.
\end{definition}

\begin{definition}[State space]\label{def:statespace}
A \emph{state space} is a tuple $\mathsf{SS} = (\St, \Trans, \state_0)$
where $\state_0 \in \St$ is the initial state.  The \emph{reachable
fragment} is
\[
  \reach(\mathsf{SS}) \;=\;
  \{\,\state \in \St \mid \state_0 \xrightarrow{*} \state\,\}
\]
where $\xrightarrow{*}$ is the reflexive-transitive closure of
$\Trans$.
\end{definition}

\subsection{Implementation: \texttt{StateSpace} and \texttt{StateExplorer}}

The Python class \texttt{StateSpace} (in
\texttt{generation/state\_space/models.py}) realises
Definition~\ref{def:statespace}.  Its core data structure is a
directed graph backed by an adjacency list; nodes are
\texttt{SemanticState} objects keyed by a SHA-256 hash of their
$(\alpha, \Omega, \mathcal{T})$ triple.

\begin{lstlisting}[style=jugeo-python, caption={Core \texttt{SemanticState} fields (simplified).}]
@dataclass
class SemanticState:
    state_id:    str                        # SHA-256 hash
    assignments: dict[str, str]             # patch -> section
    open_obligs: frozenset[str]             # open obligation ids
    treaties:    frozenset[str]             # active treaty ids
    is_terminal: bool = False

    def apply_transition(self, t: StateTransition) -> SemanticState:
        new_a = dict(self.assignments)
        if t.kind == TransitionKind.PROPOSE:
            new_a[t.patch] = t.section
        elif t.kind == TransitionKind.RETRACT:
            del new_a[t.patch]
        # ... other kinds ...
        return SemanticState(
            state_id    = _hash(new_a, ...),
            assignments = new_a,
            ...
        )
\end{lstlisting}

The \texttt{StateExplorer} class holds a reference to a
\texttt{StateSpace}, a \texttt{SearchStrategy} enum value, a
\emph{visited set} $\visited$, and a \emph{frontier} $\frontier$.  At
each step it:
\begin{enumerate}[label=(\arabic*)]
  \item Selects the next state $\state$ from $\frontier$ (according to
    the current strategy).
  \item Generates the successors of $\state$ via all enabled moves.
  \item Filters out states already in $\visited$.
  \item Adds new successors to $\frontier$ and $\visited$.
  \item Emits $\state$ to the \texttt{LocalConstructor}.
\end{enumerate}

\begin{lstlisting}[style=jugeo-python, caption={\texttt{StateExplorer.step} (simplified).}]
def step(self) -> Optional[SemanticState]:
    if not self.frontier:
        return None
    state = self._strategy_pop()        # BFS: deque.popleft(); DFS: stack.pop()
    for move in self.space.enabled_moves(state):
        succ = state.apply_transition(move)
        if succ.state_id not in self.visited:
            self.visited.add(succ.state_id)
            self._strategy_push(succ)
    self.local_constructor.consume(state)
    return state
\end{lstlisting}

The visited set is a Python \texttt{set} of state-id strings, ensuring
$O(1)$ membership tests.  Each unique state is therefore consumed by
\texttt{LocalConstructor} exactly once, eliminating redundant section
reconstruction.

\subsection{Reachability and the finite-state assumption}

\begin{definition}[Finite-state program]\label{def:finite}
A program $P$ is \emph{finite-state} if the patch cover $\patchset$ is
finite and each patch admits finitely many candidate sections.
Equivalently, $|\reach(\mathsf{SS})| < \infty$.
\end{definition}

All programs that \JuGeo{} currently handles are finite-state: the
patch cover is determined by a finite CFG, and sections are drawn from
a finite SMT-guided candidate pool.  This assumption is made explicit
in \S\ref{sec:theorem} and formalised in \LEAN{}.

% =====================================================================
\section{Exploration Strategies}\label{sec:strategies}
% =====================================================================

\subsection{The \texttt{SearchStrategy} enum}

The \texttt{SearchStrategy} enum in
\texttt{generation/state\_space/search\_strategies.py} provides four
values:

\begin{center}
\begin{tabular}{llp{7.5cm}}
\toprule
\textbf{Variant} & \textbf{Frontier structure} & \textbf{Characteristics} \\
\midrule
\texttt{BFS}   & FIFO queue      & Complete; finds min-depth goal; $O(b^d)$ memory \\
\texttt{DFS}   & LIFO stack      & Complete iff acyclic; $O(bd)$ memory \\
\texttt{BMC}   & FIFO + depth bound & Incomplete; bounded completeness guarantee \\
\texttt{HEUR}  & Priority queue  & Incomplete in general; guided by $\heur$ \\
\bottomrule
\end{tabular}
\end{center}

Here $b$ is the branching factor (number of enabled moves per state)
and $d$ is the depth of the shallowest goal state.

\subsection{Breadth-first exploration (\BFS{})}

\BFS{} maintains $\frontier$ as a FIFO queue.  It guarantees that every
state at depth $k$ is processed before any state at depth $k+1$.
Because the visited set prevents re-enqueuing, each state is dequeued
at most once.

\begin{proposition}[\BFS{} terminates on finite-state spaces]
\label{prop:bfs-terminates}
If $|\reach(\mathsf{SS})| = N < \infty$, then \BFS{} terminates after
exactly $N$ dequeue operations.
\end{proposition}

\begin{proof}
By induction on $d$: at each depth level $k$, all states reachable in
exactly $k$ steps from $\state_0$ are distinct (by the hash-keying
invariant) and each is enqueued exactly once.  Since $N$ is finite,
the frontier eventually empties. \qed
\end{proof}

\subsection{Depth-first exploration (\DFS{})}

\DFS{} maintains $\frontier$ as a LIFO stack.  In the presence of
cycles (which can arise from \textsc{Retract} followed by
\textsc{Propose} on the same patch), \DFS{} without a visited set
would loop.  With the visited set, \DFS{} also terminates on finite
spaces (same argument as \BFS{}), but visits states in a different
order that may bias toward deep, narrow paths.

\begin{remark}
\DFS{} is preferable when the first valid goal state is more important
than the minimum-depth one.  In practice, \DFS{} finds a complete
local-section assignment faster than \BFS{} when the goal state is
deep but reachable via a single path.
\end{remark}

\subsection{Bounded model checking (\BMC{})}

\BMC{} limits exploration to states reachable within depth $\depth$:

\begin{definition}[\BMC{} strategy]\label{def:bmc}
Given a depth bound $\depth \in \mathbb{N}$, \BMC{} uses a FIFO queue
augmented with depth annotations.  A state $\state$ at depth $k \ge
\depth$ is not expanded: its successors are not enqueued.  If a goal
state is reachable within depth $\depth$, \BMC{} finds it; otherwise
it declares a \emph{bounded non-termination}.
\end{definition}

The key parameter $\depth$ trades off coverage against time.
Choosing $\depth = \infty$ recovers \BFS{}.

\begin{proposition}[\BMC{} partial completeness]\label{prop:bmc}
Let $p \in [0,1]$ be the probability that a uniformly random
successor of any state is a new patch coverage.  Then \BMC{} with
bound $\depth$ covers a fraction
\[
  \covfrac_{\BMC}(\depth) \;\ge\; 1 - (1-p)^{\depth}
\]
of the total coordinate patches in expectation.
\end{proposition}

\begin{proof}
Each exploration step covers a new patch with probability at least
$p$.  After $\depth$ steps, the probability that a particular patch
remains uncovered is at most $(1-p)^{\depth}$.  Summing over all
patches and using linearity of expectation gives the bound.  \qed
\end{proof}

\subsection{Heuristic search (\HEUR{})}

\HEUR{} replaces the FIFO or LIFO discipline with a priority queue
ordered by a \emph{heuristic function} $\heur\colon \St \to \mathbb{R}$.
The default heuristic in \JuGeo{} is the \emph{coverage fraction}:
\[
  \heur(\state) \;=\;
  -\frac{|\mathrm{dom}(\alpha)|}{|\patchset|}
\]
(negated because Python's \texttt{heapq} is a min-heap).  States with
more assigned patches are explored first.  When $\heur$ is an
admissible heuristic (\ie{} never overestimates the remaining cost),
\HEUR{} reduces to A$^*$ and is optimal.

\begin{lstlisting}[style=jugeo-python, caption={Priority-queue pop for \HEUR{}.}]
class BestFirstStrategy(SearchStrategy):
    def pop(self, frontier: list) -> SemanticState:
        _, state = heapq.heappop(frontier)
        return state

    def push(self, frontier: list, state: SemanticState) -> None:
        priority = -len(state.assignments) / max(1, self.num_patches)
        heapq.heappush(frontier, (priority, state))
\end{lstlisting}

% =====================================================================
\section{Local Construction}\label{sec:local}
% =====================================================================

\subsection{\texttt{LocalConstructor}}

The \texttt{LocalConstructor} class (in
\texttt{generation/local\_construction/}) accepts a semantic state
$\state$ and produces a \emph{local section} $\LocalSec_{\state}$ at
the coordinate encoded in $\state.assignments$.

\begin{definition}[Local section from state]\label{def:locsec}
Given a semantic state $\state = (\alpha, \Omega, \mathcal{T})$ and a
coordinate $\coord \in \mathrm{dom}(\alpha)$, the \emph{local section
at $\coord$ induced by $\state$} is the restriction
\[
  \LocalSec_{\coord}^{\state} \;=\; \alpha(\coord)
\]
together with the sub-bundle of $\Omega$ and $\mathcal{T}$ scoped to
$\coord$.
\end{definition}

The construction process uses three sub-algorithms:

\begin{enumerate}
  \item \textbf{Candidate proposal} (\texttt{LoopStatus.RUNNING}):
    the constructor generates candidate sections using SMT-guided
    synthesis, type-directed enumeration, or LLM suggestion.

  \item \textbf{Interface discipline check}
    (\texttt{InterfaceDiscipline}): the candidate is tested against the
    boundary constraints $\partial\coord$ inherited from adjacent
    patches.  Two sections are \emph{compatible} iff their interface
    disciplines $D_u \sqcap D_v$ are jointly satisfiable.

  \item \textbf{Acceptance} (\texttt{LoopStatus.SUCCEEDED}): if the
    candidate passes all discipline checks and closes the scoped
    obligations, it is recorded as $\LocalSec_{\coord}^{\state}$ and
    the loop status transitions to \texttt{SUCCEEDED}.
\end{enumerate}

\begin{lstlisting}[style=jugeo-python, caption={\texttt{LocalConstructor.consume} (simplified).}]
def consume(self, state: SemanticState) -> ConstructionStatus:
    for coord, section_id in state.assignments.items():
        loop = self._get_or_create_loop(coord)
        if loop.status == LoopStatus.SUCCEEDED:
            continue                           # already built
        candidate = self._synthesise(coord, section_id, state)
        if self._check_interface(candidate, coord):
            loop.accept(candidate)
            self._status[coord] = ConstructionStatus.COMPLETE
        else:
            loop.reject(candidate)
            self._status[coord] = ConstructionStatus.IN_PROGRESS
    return self._aggregate_status()
\end{lstlisting}

\subsection{\texttt{ConstructionStatus} and state machine}

Each coordinate's construction follows the status machine:
\[
  \texttt{PENDING}
  \;\xrightarrow{\text{init}}\;
  \texttt{IN\_PROGRESS}
  \;\xrightarrow{\text{accept}}\;
  \texttt{COMPLETE}
  \;\xrightarrow{}\;
  \texttt{VERIFIED}
\]
with an additional \texttt{FAILED} sink reachable from
\texttt{IN\_PROGRESS} when the budget is exhausted.  The orchestrator
tracks this status for all coordinates and reports aggregate progress
to the explorer.

% =====================================================================
\section{Construction Orchestration}\label{sec:orchestration}
% =====================================================================

\subsection{Multi-coordinate management}

When the state explorer visits a state $\state$, multiple coordinates
in $\mathrm{dom}(\alpha)$ may be ready for local section construction.
The \texttt{ConstructionOrchestrator} manages this concurrently:

\begin{definition}[Orchestrator]\label{def:orch}
The \emph{construction orchestrator} $\Orch$ is a function
\[
  \Orch\colon
  \St \times \LocSecs^{<\omega}
  \;\longrightarrow\;
  \LocSecs^{<\omega} \times \CStatus
\]
that takes the current state and the already-built sections, attempts
construction for all new coordinates in $\mathrm{dom}(\alpha(\state))
\setminus \mathrm{dom}(\mathcal{L})$, and returns an updated section
pool together with a status indicator.
\end{definition}

The orchestrator ensures:
\begin{itemize}
  \item \textbf{Idempotence}: constructing a section for a coordinate
    that already has one in $\mathcal{L}$ is a no-op.
  \item \textbf{Monotonicity}: $|\mathcal{L}|$ is non-decreasing across
    all orchestrator calls.
  \item \textbf{Interface consistency}: newly added sections are
    checked for compatibility with all existing adjacent sections.
\end{itemize}

\subsection{Coordination graph}

The orchestrator maintains a \emph{coordination graph}
$G = (V, E)$ where $V = \patchset$ and $(u,v) \in E$ iff patch $u$
directly depends on patch $v$ (\eg{} a function call, a shared
variable, or a treaty boundary).  When a new section is accepted at
$u$, all neighbours $v \in N_G(u)$ are notified so they can update
their interface discipline $D_v$.

\begin{lstlisting}[style=jugeo-python, caption={Orchestrator dispatch (simplified).}]
class ConstructionOrchestrator:
    def dispatch(self, state: SemanticState) -> ConstructionStatus:
        new_coords = (
            set(state.assignments) - set(self.built_sections)
        )
        for coord in self._topological_order(new_coords):
            result = self.constructor.consume_coord(coord, state)
            if result == ConstructionStatus.COMPLETE:
                self.built_sections[coord] = result
                self._notify_neighbours(coord)
        return self._aggregate()
\end{lstlisting}

\subsection{Parallelism}

The orchestrator supports two parallelism modes:

\begin{enumerate}
  \item \textbf{Sequential}: coordinates are constructed in a
    topological order of $G$.  This ensures that interface discipline
    constraints from upstream patches are available before downstream
    patches are processed.

  \item \textbf{Parallel with barrier}: independent coordinates
    (those with no shared edges in $G$) are constructed concurrently
    using Python \texttt{asyncio} tasks.  A barrier synchronises at
    each topological level before proceeding to the next.
\end{enumerate}

In the sequential mode, the orchestrator is a special case of a
\emph{coordinated elaboration} (Paper~14, \S4): the coordination graph
defines the elaboration order.

% =====================================================================
\section{Completion Conditions}\label{sec:completion}
% =====================================================================

\subsection{The \texttt{CompletionCondition} abstraction}

Exploration terminates when a \texttt{CompletionCondition} is
satisfied.  The abstraction separates the \emph{what} (completion
criterion) from the \emph{how} (search strategy):

\begin{definition}[\texttt{CompletionCondition}]\label{def:completion}
A \emph{completion condition} is a predicate
\[
  \CCond\colon \LocSecs^{<\omega} \times \mathbb{N} \times \St
  \;\longrightarrow\; \{\mathsf{true},\mathsf{false}\}
\]
taking the current section pool $\mathcal{L}$, the number of explored
states $k$, and the current state $\state$.  The explorer halts as
soon as $\CCond(\mathcal{L}, k, \state)$ is true.
\end{definition}

The implementation provides four concrete conditions:

\begin{center}
\begin{tabular}{lp{9cm}}
\toprule
\textbf{Condition} & \textbf{Description} \\
\midrule
\texttt{ALL\_PATCHES}    & All patches in $\patchset$ have an accepted section ($\covfrac = 1$). \\
\texttt{BUDGET}          & The number of explored states $k$ exceeds a fixed budget $\budg$. \\
\texttt{COVERAGE\_FRAC}  & $\covfrac \ge \theta$ for a configurable threshold $\theta$. \\
\texttt{FRONTIER\_EMPTY} & The frontier $\frontier$ is empty (full exhaustion). \\
\bottomrule
\end{tabular}
\end{center}

Conditions can be combined conjunctively (\texttt{AND}) or
disjunctively (\texttt{OR}) to express policies such as ``stop when
all patches are covered \emph{or} the budget is exceeded.''

\subsection{Interaction with strategies}

Each strategy pair with a natural default condition:

\begin{itemize}
  \item \BFS{}: \texttt{FRONTIER\_EMPTY} (full reachability) or
    \texttt{ALL\_PATCHES} (early exit on coverage).
  \item \DFS{}: \texttt{ALL\_PATCHES} (stop at first goal state).
  \item \BMC{}: \texttt{BUDGET} with $\budg = |\reach(\mathsf{SS}, \depth)|$.
  \item \HEUR{}: \texttt{COVERAGE\_FRAC} with $\theta = 0.95$.
\end{itemize}

\subsection{Completion certificate}

When $\CCond$ fires, the orchestrator issues a \emph{completion
certificate}: a record containing the final section pool $\mathcal{L}$,
the coverage fraction $\covfrac$, the number of explored states, and
the elapsed time.  This certificate is stored in the judgment 8-tuple
as part of the evidence bundle $\evid$.

% =====================================================================
\section{Bounded Completeness Theorem}\label{sec:theorem}
% =====================================================================

We now state and prove the main theoretical result.

\begin{theorem}[Bounded Completeness]\label{thm:bounded-completeness}
Let $P$ be a finite-state program with reachable state space
$\reach(\mathsf{SS}) = \{\state_0, \dots, \state_{N-1}\}$ and patch
cover $\patchset$ of size $n$.  Let $\mathcal{L}^*$ be the union of
all local sections obtainable by any sequence of semantic moves from
$\state_0$.  Then \BFS{} with completion condition
\texttt{FRONTIER\_EMPTY} produces a section pool $\mathcal{L}$ such
that:
\begin{enumerate}[label=(\roman*)]
  \item $\mathcal{L} = \mathcal{L}^*$ \quad (completeness),
  \item Each state in $\reach(\mathsf{SS})$ is visited exactly once
    \quad (no redundancy),
  \item The total number of steps is $O(N + |\Trans|)$ \quad
    (time proportional to the reachable state space).
\end{enumerate}
\end{theorem}

\begin{proof}
We prove each part in turn.

\smallskip\noindent
\emph{(i) Completeness.}
We show that $\mathcal{L} \supseteq \mathcal{L}^*$.  Let
$\LocalSec \in \mathcal{L}^*$ be any section constructible from
$\state_0$.  Then there exists a path $\state_0 \xrightarrow{m_1}
\state_1 \xrightarrow{m_2} \cdots \xrightarrow{m_k} \state_k$ with
$\LocalSec \in \mathrm{sections}(\state_k)$.  By induction on $k$:
the base case $k=0$ holds since $\state_0$ is enqueued initially.
For the inductive step, assume $\state_{k-1}$ is enqueued.  When
$\state_{k-1}$ is dequeued, \BFS{} generates all successors, including
$\state_k$.  Since $\state_k \notin \visited$ (it has not been reached
by any shorter path; if it had, its section would already be in
$\mathcal{L}$), it is enqueued and eventually dequeued.  When
$\state_k$ is processed, $\mathrm{LocalConstructor}$ adds $\LocalSec$
to $\mathcal{L}$.  Hence $\mathcal{L} \supseteq \mathcal{L}^*$.

The reverse inclusion $\mathcal{L} \subseteq \mathcal{L}^*$ is
immediate: \texttt{LocalConstructor} only adds sections reachable from
$\state_0$.

\smallskip\noindent
\emph{(ii) No redundancy.}
Each state $\state \in \reach(\mathsf{SS})$ is added to $\visited$ at
most once (when first discovered).  A state in $\visited$ is never
re-enqueued.  Therefore each state is dequeued at most once.  Since
every reachable state is eventually enqueued (by part (i)), it is
dequeued exactly once.

\smallskip\noindent
\emph{(iii) Time complexity.}
Each state is dequeued once ($N$ dequeue operations).  Each dequeue
generates the successors of $\state$: the number of successors is
bounded by the out-degree of $\state$ in $\Trans$.  Summing over all
states, the total number of successor-generation operations is
$|\Trans|$.  Each visited-set membership test and insertion is $O(1)$
(hash set).  Therefore the total time is $O(N + |\Trans|)$. \qed
\end{proof}

\begin{corollary}[BFS coverage]\label{cor:bfs-coverage}
Under the conditions of Theorem~\ref{thm:bounded-completeness}, the
coverage fraction after \BFS{} terminates satisfies $\covfrac = 1$.
\end{corollary}

\begin{proof}
By part (i) of the theorem, $\mathcal{L} = \mathcal{L}^*$.  Since
$P$ is finite-state and all patches admit at least one constructible
section (by the finite-state assumption), $\mathrm{dom}(\mathcal{L})
= \patchset$, so $\covfrac = |\mathrm{dom}(\mathcal{L})| / |\patchset|
= n/n = 1$. \qed
\end{proof}

\begin{corollary}[DFS completeness under acyclicity]\label{cor:dfs}
If $\Trans$ is acyclic (the state space is a DAG), then \DFS{} with
the visited set also achieves $\covfrac = 1$.
\end{corollary}

\begin{proof}
In a DAG, \DFS{} with a visited set visits every reachable node exactly
once by the same argument as \BFS{}; the LIFO discipline does not affect
reachability. \qed
\end{proof}

\begin{remark}
When $\Trans$ has cycles, \DFS{} without a visited set may loop
indefinitely.  The \texttt{StateExplorer} \emph{always} maintains a
visited set, so Corollary~\ref{cor:dfs} applies to the implementation
even when the underlying state space has cycles.
\end{remark}

% =====================================================================
\section{Experiments}\label{sec:experiments}
% =====================================================================

\subsection{Benchmark setup}

We evaluate the three main strategies (\BFS{}, \DFS{}, \HEUR{}) on the
\numcases{}-program benchmark suite used throughout the \JuGeo{} series.
Programs range from 10 to 500 lines of Python, covering sorting
algorithms, data-structure operations, arithmetic routines, and simple
protocol implementations.  The benchmark includes the \numspec{}
specification programs, \numeq{} equivalence pairs, and \numbug{} bug
programs from Papers~1--10.

For each program we measure:
\begin{itemize}
  \item \textbf{Coverage fraction} $\covfrac$: fraction of patches with
    an accepted local section.
  \item \textbf{Explored states}: number of \texttt{SemanticState}
    objects processed.
  \item \textbf{Wall time}: end-to-end time for the explorer loop.
  \item \textbf{Round count}: number of orchestrator dispatch calls.
\end{itemize}

All experiments run on a single core of a 3.2\,GHz machine with 16\,GB
RAM.  The SMT backend is Z3~4.13.

\subsection{Strategy comparison}

\begin{table}[H]
\centering
\caption{Exploration strategy comparison on \numcases{} programs.
  ``States'' = mean explored states per program;
  ``Time'' = mean wall time;
  ``$\covfrac$'' = mean coverage fraction;
  ``Rounds'' = mean orchestrator dispatch rounds.}
\label{tab:strategies}
\begin{tabular}{lccccc}
\toprule
\textbf{Strategy} & \textbf{States} & \textbf{Time (ms)} &
  \textbf{$\covfrac$} & \textbf{Rounds} & \textbf{Completeness} \\
\midrule
\BFS{}             & \ppXXXIBfsStates   & \ppXXXIBfsTimeMs  & \ppXXXIBfsCoverage & \ppXXXIBfsRounds & Full \\
\DFS{}             & \ppXXXIDfsStates   & \ppXXXIDfsTimeMs  & \ppXXXIDfsCoverage & \ppXXXIDfsRounds & Full \\
\BMC{} ($\depth=5$)& \ppXXXIBmcStates   & \ppXXXIBmcTimeMs   & \ppXXXIBmcCoverage & \ppXXXIBmcRounds & Bounded \\
\HEUR{}            & \ppXXXIHeurStates   & \ppXXXIHeurTimeMs   & \ppXXXIHeurCoverage & \ppXXXIHeurRounds & Full \\
\bottomrule
\end{tabular}
\end{table}

\BFS{} and \DFS{} achieve full coverage ($\covfrac = 1$) on all programs (as
guaranteed by the theorem).  \HEUR{} matches full coverage while
exploring substantially fewer states than \BFS{} requires.
\BMC{} at bounded depth may miss some patches; raising the
bound increases coverage toward full.

\subsection{Round efficiency}

\DFS{} requires more orchestrator rounds on average than \BFS{}.
This is because \DFS{} visits deep states early, triggering
interface-discipline negotiation with ancestors before those ancestors
are fully constructed.  The orchestrator must defer and retry these
negotiations, inflating the round count.  \BFS{} processes ancestors
before descendants, so negotiations succeed on the first attempt.

\subsection{Scaling with state-space size}

\begin{figure}[H]
\centering
\begin{tabular}{cccc}
\toprule
\textbf{Program size (lines)} & \textbf{$|\reach|$} &
  \textbf{BFS time (ms)} & \textbf{DFS time (ms)} \\
\midrule
$\le 50$  & $\le 120$  & \ppXXXISmallBfsMs & \ppXXXISmallDfsMs \\
$51$--$100$ & $121$--$400$ & \ppXXXIMedBfsMs & \ppXXXIMedDfsMs \\
$101$--$250$ & $401$--$1500$ & \ppXXXILargeBfsMs & \ppXXXILargeDfsMs \\
\emph{No xlarge programs in benchmark} & --- & --- & --- \\
\bottomrule
\end{tabular}
\caption{BFS and DFS wall time vs.\ program size.}
\label{tab:scaling}
\end{figure}

Times scale near-linearly with $|\reach|$, consistent with the $O(N +
|\Trans|)$ bound.  The slight super-linear component at larger sizes
is attributable to hash-collision overhead in the visited set as
$|\visited|$ grows beyond the L3 cache.

\subsection{Completion condition sensitivity}

We compare two completion conditions for \HEUR{}:
\texttt{ALL\_PATCHES} (stop when $\covfrac = 1$) vs.\
\texttt{COVERAGE\_FRAC} with a high coverage threshold.  On programs where the
last patches require many high-depth states,
\texttt{COVERAGE\_FRAC} reduces the state budget significantly
with a small coverage cost.  For
verification-generation applications where high but not perfect
coverage is acceptable, this trade-off is favourable.

% =====================================================================
\section{Conclusion}\label{sec:conclusion}
% =====================================================================

We have developed a principled state space exploration framework for
exhaustive local section generation in \JuGeo{}.  The key contributions
are:

\begin{enumerate}
  \item A formal \textbf{state space model} that captures the
    generation process as a directed graph of semantic states, with a
    precise notion of reachability (\S\ref{sec:statespace}).

  \item A \textbf{pluggable strategy} design---\BFS{}, \DFS{}, \BMC{},
    \HEUR{}---that separates the frontier discipline from the
    termination criterion (\S\ref{sec:strategies}).

  \item A \textbf{local construction integration} showing how each
    visited state feeds the \texttt{LocalConstructor} and
    \texttt{ConstructionOrchestrator} (\S\S\ref{sec:local}--\ref{sec:orchestration}).

  \item A \textbf{Bounded Completeness Theorem} establishing that BFS
    achieves $\covfrac = 1$ for finite-state programs in $O(N +
    |\Trans|)$ time, with no redundant state reconstruction
    (\S\ref{sec:theorem}).
\end{enumerate}

The framework closes the coverage gap identified in Papers~1--30 and
provides the theoretical foundation for exhaustive evidence generation
in the \JuGeo{} system.

\paragraph{Future work.}
Three directions are most promising:
(a)~\emph{Incremental exploration}: when a program changes, re-explore
only the affected sub-graph of the state space.
(b)~\emph{Compositional state spaces}: compose state spaces for
sub-programs using the operadic machinery of Paper~14.
(c)~\emph{Probabilistic strategies}: assign weights to transitions
derived from coverage likelihood, enabling a principled
\emph{probabilistic completeness} result analogous to
Proposition~\ref{prop:bmc} but adaptive.

\begin{thebibliography}{99}

\bibitem{jugeo01}
H.~Young.
\newblock Semantic sites for program verification.
\newblock {\em Judgment Geometry, Paper~1}, 2024.

\bibitem{jugeo02}
H.~Young.
\newblock Judgment algebra: An 8-tuple framework.
\newblock {\em Judgment Geometry, Paper~2}, 2024.

\bibitem{jugeo03}
H.~Young.
\newblock Sheaf descent and obstruction classes.
\newblock {\em Judgment Geometry, Paper~3}, 2024.

\bibitem{jugeo06}
H.~Young.
\newblock Semantic moves for \textsc{JuGeo}.
\newblock {\em Judgment Geometry, Paper~6}, 2024.

\bibitem{jugeo14}
H.~Young.
\newblock Operadic composition of verification strategies.
\newblock {\em Judgment Geometry, Paper~14}, 2024.

\bibitem{jugeo30}
H.~Young.
\newblock Motivic measures for cross-language verification transfer.
\newblock {\em Judgment Geometry, Paper~30}, 2025.

\bibitem{biere2003bounded}
A.~Biere, A.~Cimatti, E.~M.~Clarke, O.~Strichman, and Y.~Zhu.
\newblock Bounded model checking.
\newblock {\em Advances in Computers}, 58:117--148, 2003.

\bibitem{clarke1999model}
E.~M.~Clarke, O.~Grumberg, and D.~Peled.
\newblock {\em Model Checking}.
\newblock MIT Press, 1999.

\bibitem{hart1968formal}
P.~E.~Hart, N.~J.~Nilsson, and B.~Raphael.
\newblock A formal basis for the heuristic determination of minimum cost paths.
\newblock {\em IEEE Trans.\ Syst.\ Sci.\ Cybernetics}, 4(2):100--107, 1968.

\bibitem{moura2021lean4}
L.~de Moura and S.~Ullrich.
\newblock The {Lean~4} theorem prover and programming language.
\newblock {\em Proc.\ CADE 2021}, LNAI 12699.

\bibitem{king1976symbolic}
J.~C.~King.
\newblock Symbolic execution and program testing.
\newblock {\em Commun.\ ACM}, 19(7):385--394, 1976.

\bibitem{cousot1977abstract}
P.~Cousot and R.~Cousot.
\newblock Abstract interpretation: A unified lattice model for static analysis.
\newblock {\em Proc.\ POPL 1977}.

\bibitem{z3}
L.~M.~de~Moura and N.~Bj{\o}rner.
\newblock {Z3}: An efficient {SMT} solver.
\newblock {\em Proc.\ TACAS 2008}, LNCS~4963.

\end{thebibliography}

\end{document}
